\section{Second-Order Linear Types}
\label{sec:second}

This section studies the construction of Section~\ref{sec:recipe} at
second-order types (types of programs whose arguments are functions),
where it returns a strict subset of the probabilistic coherence space. A
formula has a \emph{gap} if its calculus is strictly smaller than its
probabilistic coherence space. Section~\ref{sec:second-seq} treats the
type $T$ of a program with two functional arguments,
Section~\ref{sec:second-other} extends the strict inclusion to every
formula with a negative occurrence of the core of $T$, the formula whose
negation is $T$, and Sections~\ref{sec:second-shannon}
and~\ref{sec:second-perfect} relate the gap to coherence spaces and to
graph polytopes and characterize it at second order. By computation, not formalized, the construction and the
probabilistic coherence space agree on $(\mathrm{Bool} \multimap
\mathrm{Bool}) \multimap \mathrm{Bool}$ and $((\mathrm{Bool} \multimap
\mathrm{Bool}) \multimap \mathrm{Bool}) \multimap \mathrm{Bool}$, the two
iterated implications into $\mathrm{Bool}$ we checked.

\subsection{Randomized Sequential Programs}
\label{sec:second-seq}

At second order, let
$T = ((X \multimap Y) \otimes (X' \multimap Y'))^{\perp}$, the type of a
program that receives $f : X \multimap Y$ and $g : X' \multimap Y'$,
uses each once, and halts or diverges. Its negation
$(X \multimap Y) \otimes (X' \multimap Y')$ is the \emph{core} of $T$. We
take $T$ to be the set of
$t : (X \times Y) \times (X' \times Y') \to \ENNReal$ with
$\sum_{p,q} y_p\, t(p,q)\, z_q \le 1$ for all $y$ in $X \multimap Y$ and
$z$ in $X' \multimap Y'$, the orthogonal of the products $y \otimes z$.
For $X = Y = X' = Y' = \mathrm{Bool}$, Lean checks that $T$ satisfies
the three PCS conditions\leanref{ispcs-lin}. A \emph{randomized sequential program} queries
one argument first. Querying $f$ first, it chooses the input $a$ with
probability $\mu(a)$, after the answer $b$ chooses the input $c$ to $g$
with probability $\kappa(a,b,c)$, and after the answer $d$ halts with
probability $\rho(a,b,c,d)$. Its value at $((a,b),(c,d))$ is
$\mu(a)\kappa(a,b,c)\rho(a,b,c,d)$ plus the analogous term for
programs that query $g$ first, with $\sum \mu \le 1$ over both orders,
$\sum_c \kappa(a,b,c) \le 1$, and $\rho \le 1$. It is
\emph{deterministic} if its values are in $\{0,1\}$.

\begin{theorem}[second order\leanref{sharp-td-iff}\leanref{recipe-eq-sequentialg}]
\label{thm:second-order}
For nonempty finite $X, X'$ and finite $Y, Y'$:
(i) A $\{0,1\}$-valued $z$ lies in $T$ iff its support lies in a tree
$\{((a, b), (c(b), d)) \mid b \in Y, d \in Y'\}$ for some $a \in X$ and
$c : Y \to X'$, or in a tree $\{((a(d), b), (c, d))\}$ that queries $g$
first, and these are exactly the deterministic sequential
programs\leanref{sharp-mem-sequentialg}.
(ii) $\mathrm{Mix}$ of the $\{0,1\}$-valued elements of $T$ is exactly the
set of randomized sequential programs.
\end{theorem}

\noindent
Part (ii) contains both directions of Kuhn's theorem for this
type~\cite{kuhn1953}: every randomized sequential program is a mixture of
deterministic ones\leanref{sequentialg-subset-mix}, and every mixture of
randomized sequential programs is a randomized sequential
program\leanref{mix-sequentialg}. The second direction conditions the
choices at each node on the history, and nodes reached with probability
$0$ are given weight $0$. For $X = Y = X' = Y' = \mathrm{Bool}$ the same statements are also
proved directly\leanref{kuhn-bool}.

\paragraph{A language for $T$.}
Randomized sequential programs are the denotations of a small
language\leanref{proglang}. It has four syntactic stages: no function
called, $f$ called, $g$ called, and both called. Each stage has
divergence and a finite probabilistic choice
$\mathsf{choose}(w_1 P_1, \dots, w_k P_k)$ with $\sum_i w_i \le 1$, the
missing mass diverging. The first stage has $\mathsf{call}_f\,a\,k$ and
$\mathsf{call}_g\,c\,k$ with a continuation $k$ for each answer, the
middle stages call the remaining function, and the last stage halts. The
stages make each function called exactly once. A choice denotes the
weighted sum of the denotations of its branches, and $\mathsf{call}_f\,a\,k$
denotes $((a',b),(c,d)) \mapsto [a' = a]\,\llbracket k\,b\rrbracket(c,d)$.

\begin{theorem}[soundness and definability\leanref{proglang-sound}\leanref{proglang-definable}]
\label{thm:proglang}
For nonempty finite $X, X'$ and finite $Y, Y'$, every program denotes a
randomized sequential program, hence an element of $\mathrm{Mix}$ of the
$\{0,1\}$-valued elements of $T$. Every $\{0,1\}$-valued element of $T$
is the denotation of a program without choice. So $\mathrm{Mix}$ of the
$\{0,1\}$-valued elements of $T$ is exactly the set of denotations of
programs\leanref{proglang-eq}.
\end{theorem}

\noindent The
language is not a full programming language: it has no recursion and no
higher-order arguments beyond $f$ and $g$. Relating it to probabilistic
PCF at this type is not done here.

The calculus is strictly smaller than the PCS. Call two distinct web
points of $T$ \emph{exclusive} if a single pair of deterministic partial
programs $(f, g)$ produces both, that is, if their first components
cohere in $X \multimap Y$ and their second components cohere in
$X' \multimap Y'$. The \emph{exclusivity graph} of $T$ has the web points
as vertices and the exclusive pairs as edges.

\begin{theorem}[strict inclusion\leanref{pente-not-mix}\leanref{pent-bounds}]
\label{thm:pentagon}
Let $X, Y, X', Y'$ have at least two elements each. There are five web
points of $T$ such that the function with value $\tfrac12$ on them and
$0$ elsewhere lies in $T$ and is not in $\mathrm{Mix}$ of the
$\{0,1\}$-valued elements of $T$. For $X = Y = X' = Y' = \mathrm{Bool}$, the total weight on the five
points is at most $2$ for every randomized sequential program, $2$ is
attained, and it is at most $\tfrac52$ on $T$, attained by the function
above.
\end{theorem}

\noindent
Over $\mathrm{Bool}$, writing $(a,b)$ for the observation $f(a) = b$, the
five points are $((t,t),(t,t))$, $((t,f),(f,f))$, $((f,t),(t,f))$,
$((f,t),(f,t))$, and $((f,f),(f,f))$. They induce a $5$-cycle in the
exclusivity graph. Lean checks the two facts the proof uses, that the
induced subgraph has no triangle and no three pairwise non-adjacent
vertices\leanref{pent-cycle}, and
embeds the five points into larger types by injections from
$\mathrm{Bool}$.

The witness is partial: for some pairs of total functions it halts with
probability less than $1$. At $T$ this is necessary. An element of $T$
that halts with probability $1$ for every pair of total functions is a
two-party classical process without predefined causal order (a process
in which the order of the two calls may depend on the inputs), and every
such process is a mixture of processes in a fixed
order~\cite{baumeler-wolf2016}. For $X = Y = X' = Y' = \mathrm{Bool}$ we
checked by vertex enumeration, not formalized, that these elements form a
polytope whose $12$ vertices are $\{0,1\}$-valued. With three functional
arguments the gap remains among the elements that halt with probability
$1$.

\begin{proposition}[a gap among total programs\leanref{gap-t3}]
\label{prop:t3}
Let $T_3 = ((\mathbf{2} \multimap \mathbf{2}) \otimes (\mathbf{2}
\multimap \mathbf{2}) \otimes (\mathbf{2} \multimap \mathbf{2}))^{\perp}$.
Some element of $P(T_3)$ halts with probability $1$ on every triple of
total functions and is not in the calculus of $T_3$.
\end{proposition}

\noindent
The witness has value $1$ at one point of the web and $\tfrac12$ at
fourteen others. Every $\{0,1\}$-valued element of
$P(\mathbf{2} \multimap \mathbf{2})$ and of
$P((\mathbf{2} \multimap \mathbf{2}) \otimes (\mathbf{2} \multimap
\mathbf{2}))$ lies below the graph of a total function or a product of
two, so by the description of $P((X \otimes Y)^{\perp})$ in
Section~\ref{sec:second-perfect} membership reduces to the $64$ triples
of total functions, on each of which the witness has mass exactly $1$.
Five of its points of value $\tfrac12$ are pairwise exclusive along a
$5$-cycle, so, as for Theorem~\ref{thm:pentagon}, every mixture of
$\{0,1\}$-valued elements has weight at most $2$ on them, against
$\tfrac52$. The witness is an instance of the three-party processes of
Baumeler and Wolf that are not mixtures of deterministic
ones~\cite{baumeler-wolf2016}.

The semantics of Definition~\ref{def:mall} computes $T$:
$P(((\mathbf{n} \multimap \mathbf{m}) \otimes (\mathbf{n'} \multimap
\mathbf{m'}))^{\perp})$ is $T$ for $X = \{1,\dots,n\}$, $Y =
\{1,\dots,m\}$, $X' = \{1,\dots,n'\}$, $Y' = \{1,\dots,m'\}$\leanref{p-second}.
So Theorem~\ref{thm:second-order}(ii) and the strict inclusion of
Theorem~\ref{thm:pentagon} are statements about these
formulas\leanref{recipe-formula}.

\paragraph{Related work.}
Full abstraction (the model identifies exactly the programs that no
context distinguishes) of probabilistic coherence spaces for probabilistic
PCF~\cite{ehrhard-pagani-tasson2018} is proved from the analyticity of
the morphisms (their power-series form), not from definability of every
point, and in the concurrent game model of
Castellan, Clairambault, Paquet and Winskel~\cite{castellan2018}
definability fails until strategies are restricted to sequential
probabilistic innocent ones (strategies that call one argument at a time
and depend only on the history they observe). Kuhn's
theorem~\cite{kuhn1953} is the classical form of
Theorem~\ref{thm:second-order}(ii) for games with perfect recall (each
player remembers its own past moves and observations). The tensor of
probabilistic coherence spaces has been related to no-signalling
composition (composition in which neither component's outputs depend on
the other's inputs)~\cite{simmons-kissinger2022}. What
is added here is an explicit point outside the calculus at a type without
exponentials, its transfer to every formula with a negative occurrence of
the core (Theorem~\ref{thm:occ-gap}), the description of the calculus and
of the space as $\mathrm{STAB}$ and $\mathrm{QSTAB}$ of one graph
(Theorem~\ref{thm:qstab}), the exact condition at second order
(Theorem~\ref{thm:second-exact}), and the mechanization.

\subsection{Other Formulas}
\label{sec:second-other}

The strict inclusion transfers along a simple kind of embedding of webs.

\begin{definition}[coordinate retract\leanref{retract}]
\label{def:retract}
For formulas $D$ and $A$ of Definition~\ref{def:mall}, a \emph{coordinate
retract} of $D$ into $A$ is an injective map $j : |D| \to |A|$ such that
the pushforward $j_{*}$, with $(j_{*}x)(j(d)) = x_d$ and value $0$ outside
the image of $j$, maps $P(D)$ into $P(A)$, and the pullback
$z \mapsto z \circ j$ maps $P(A)$ into $P(D)$.
\end{definition}

\begin{theorem}[transfer\leanref{retract-gap}\leanref{retract-ops}]
\label{thm:retract}
(i) If $D$ is a coordinate retract of $A$ and the calculus of $D$ is
strictly smaller than $P(D)$, then the calculus of $A$ is strictly
smaller than $P(A)$.
(ii) Coordinate retracts compose. If $D$ is a coordinate retract of $A$,
then $D^{\perp}$ is a coordinate retract of $A^{\perp}$. The formula $D$
is a coordinate retract of $D \mathbin{\&} B$, $B \mathbin{\&} D$, and,
if $|B|$ is nonempty, of $D \otimes B$ and $B \otimes D$.
\end{theorem}

\noindent
For (i), take $x \in P(D)$ outside the calculus. If $j_{*}x$ were a
mixture of $\{0,1\}$-valued elements of $P(A)$, pulling the mixture back
along $j$ would write $x = (j_{*}x) \circ j$ as a mixture of
$\{0,1\}$-valued elements of $P(D)$. For negation, the pairing satisfies
$\sum_a (j_{*}u)_a z_a = \sum_d u_d\, z_{j(d)}$, which exchanges the two
conditions. For $D \otimes B$, the map is $j(d) = (d, b)$ for a point $b$
of $B$, so $j_{*}x = x \otimes \delta_b$, and the pullback condition uses
that every element of $P(B)$ is at most $1$ at $b$
(Theorem~\ref{thm:mall}(ii)).

An occurrence of a formula $D$ in $A$ is \emph{positive} if it lies under
an even number of negations and \emph{negative} otherwise, the derived
connectives $\parr$, $\oplus$, and $\multimap$ being unfolded. So the left
argument of $\multimap$ is a negative position.

\begin{theorem}[the gap at other formulas\leanref{gap-of-occ}]
\label{thm:occ-gap}
Let $n, m, n', m' \ge 2$, and let $A$ be a formula all of whose data types
are nonempty. If $(\mathbf{n} \multimap \mathbf{m}) \otimes
(\mathbf{n'} \multimap \mathbf{m'})$ occurs negatively in $A$, the
calculus of $A$ is strictly smaller than $P(A)$.
\end{theorem}

\noindent
A positive occurrence of $D$ gives a coordinate retract of $D$ into $A$,
and a negative one a coordinate retract of $D^{\perp}$, by
Theorem~\ref{thm:retract}(ii) and induction on $A$, using that $D$ is a
coordinate retract of $D^{\perp\perp}$. Theorem~\ref{thm:retract}(i) and
Theorem~\ref{thm:pentagon} then give the claim. Examples are $T$ itself,
the type $((\mathbf{n} \multimap \mathbf{m}) \otimes (\mathbf{n'}
\multimap \mathbf{m'})) \multimap B$ of a program that receives two
functions and returns a value of type $B$, for every $B$ with nonempty data
types (Lean checks the instance $B = \mathbf{2}$), and every formula
containing one of these positively. A
positive occurrence of the core does not suffice:
Section~\ref{sec:second-perfect} shows that at the core itself the
calculus equals the probabilistic coherence space.

The occurrence condition of Theorem~\ref{thm:occ-gap} is syntactic. The
following version depends only on
the coherence graphs of Definition~\ref{def:mall}. An \emph{induced $4$-cycle} of a formula $X$ is a
quadruple of points $a, b, c, d$ with $a \coh b \coh c \coh d \coh a$
and with $a, c$ and $b, d$ incoherent.

\begin{theorem}[four-cycles\leanref{retract-tens}\leanref{c4-retract}\leanref{gap-of-c4}]
\label{thm:c4}
If $X$ and $Y$ both have an induced $4$-cycle, the calculus of
$(X \otimes Y)^{\perp}$ is strictly smaller than its probabilistic
coherence space, and so is the calculus of every formula over nonempty
data types in which $X \otimes Y$ occurs negatively.
\end{theorem}

\noindent
The four points form a coordinate retract of $\mathbf{2} \multimap
\mathbf{2}$, with input $0$ sent to $a, c$ and input $1$ to $b, d$. The
pushforward lands in $P(X)$ because the coherent pairs of the cycle lie in
$P(X)$, and the pullback lands in $P(\mathbf{2} \multimap \mathbf{2})$
because the incoherent pairs sum to at most $1$
(Theorem~\ref{thm:mall}(ii)). Coordinate retracts pass through $\otimes$,
so Theorems~\ref{thm:pentagon} and~\ref{thm:retract} give the claim. The
coherence graph of $\mathbf{n} \multimap \mathbf{m}$ has an induced
$4$-cycle iff $n, m \ge 2$, so Theorem~\ref{thm:occ-gap} is the case of a
literal occurrence of the core. The curried type $(\mathbf{2} \multimap
\mathbf{2}) \multimap ((\mathbf{2} \multimap \mathbf{2}) \multimap
\mathbf{2})$, which Theorem~\ref{thm:occ-gap} does not cover, is an
instance\leanref{gap-curried}.

\subsection{A Clique Outside the PCS}
\label{sec:second-shannon}

The converse of Theorem~\ref{thm:mall}(iii) fails. Let
$T = ((\mathbf{2} \multimap \mathbf{2}) \otimes (\mathbf{2} \multimap
\mathbf{2}))^{\perp}$. On the five points of Theorem~\ref{thm:pentagon},
two distinct points cohere in $T$ iff they are not exclusive. Order them
as $c_0, \dots, c_4$ along the $5$-cycle of the exclusivity graph. The
coherence graph on them is the complement of this cycle, again a
$5$-cycle. Two distinct points of $T \otimes T$ are incoherent iff they
are incoherent in one component. A clique of $(T \otimes T)^{\perp}$
contained in $\{c_0, \dots, c_4\}^2$ is therefore a stable set of the
strong product of the two coherence graphs (the graph on pairs in which
two distinct pairs are adjacent iff each coordinate is equal or
adjacent), a copy of
$C_5 \boxtimes C_5$. The set $\{(c_i, c_{2i}) \mid i \in \mathbb{Z}/5\}$
is such a stable set, of size $5$. It is Shannon's stable set, which
shows that the Shannon capacity of $C_5$ (the growth rate of the largest
stable sets of its strong powers) is at least
$\sqrt5$~\cite{lovasz1979}.

\begin{theorem}[a clique outside the PCS\leanref{shannon-clique}\leanref{shannon-not-mem}]
\label{thm:shannon}
Shannon's set is a clique of the coherence space $(T \otimes T)^{\perp}$,
and its indicator is not an element of $P((T \otimes T)^{\perp})$.
\end{theorem}

\noindent
The function $z$ of Theorem~\ref{thm:pentagon}, with value $\tfrac12$ on
the five points, lies in $P(T)$, so $z \otimes z$ lies in
$P(T \otimes T)$. Its pairing with the indicator of Shannon's set is
$5 \cdot \tfrac14 = \tfrac54 > 1$. The theorem is also proved for the web
of $T$ given directly by $\mathrm{Bool}$\leanref{tensorgap-bool}. From
$(T \otimes T)^{\perp}$ on, the two natural choices of certain states
therefore differ. With the $\{0,1\}$-valued elements of the PCS, the
calculus is contained in the PCS at every formula
(Theorem~\ref{thm:mall}(iv)). With the cliques of the coherence space, it
is not.

\subsection{When the Calculus Equals the PCS}
\label{sec:second-perfect}

For a finite graph $G$, a \emph{clique} is a set of pairwise adjacent
vertices and a \emph{stable set} a set of pairwise non-adjacent
vertices. Let $\mathrm{STAB}(G)$ be $\mathrm{Mix}$ of the indicator
functions of stable sets, and $\mathrm{QSTAB}(G)$ the set of
$x : V \to \ENNReal$ with $\sum_{v \in C} x_v \le 1$ for every clique
$C$~\cite{gls1988}.

\begin{theorem}[graph polytopes\leanref{td-eq-qstab}\leanref{recipe-eq-stab}]
\label{thm:qstab}
$T = \mathrm{QSTAB}(G)$ and $\mathrm{Mix}$ of the $\{0,1\}$-valued
elements of $T$ is $\mathrm{STAB}(G)$, for $G$ the exclusivity graph of
$T$.
\end{theorem}

\begin{theorem}[odd holes and antiholes\leanref{berge-of-stab-eq-qstab}\leanref{berge-of-recipe-eq-pcs}]
\label{thm:berge}
If $\mathrm{STAB}(G) = \mathrm{QSTAB}(G)$, then $G$ contains no induced
cycle of odd length $2k+1 \ge 5$ and no induced complement of one. In
particular, if $\mathrm{Mix}$ of the $\{0,1\}$-valued elements of $T$
equals $T$, the exclusivity graph of $T$ has neither.
\end{theorem}

\noindent
A graph with neither is \emph{Berge}. A graph is \emph{perfect} if every
induced subgraph has chromatic number equal to its largest clique size.
The converse of Theorem~\ref{thm:berge} holds: $\mathrm{STAB}(G) =
\mathrm{QSTAB}(G)$ iff $G$ is perfect~\cite{chvatal1975,lovasz1972}, and
$G$ is perfect iff it is Berge, by the strong perfect graph
theorem~\cite{chudnovsky2006}. The converse is not formalized.

\paragraph{The exact condition at second order.}
A formula $X$ has the \emph{clique property} if $P(X)$ is the set of
mixtures of the indicators of the cliques of its coherence graph $G_X$.
Data types, their negations, and $\mathbf{n} \multimap \mathbf{m}$
(Theorem~\ref{thm:pcs-recipe}) have it, and so does every formula built
from data types with $\mathbin{\&}$ and $\oplus$, since $P(A \mathbin{\&}
B)$ is the product of $P(A)$ and $P(B)$ and $P(A \oplus B)$ is their convex
hull, embedded in the two summands~\cite{danos-ehrhard2011}. Lean
checks this for formulas built from $\mathbf{1}$ with $\mathbin{\&}$
and $\oplus$\leanref{recipe-code}. The coherence graphs of the latter
are the \emph{cographs}, the graphs with no induced path on four
vertices, with $\mathbin{\&}$ as join (all edges between the two parts
added) and $\oplus$ as disjoint union.

\begin{theorem}[second order\leanref{closure}\leanref{recipe-core}\leanref{second-iff}]
\label{thm:second-exact}
Let $X$ and $Y$ have the clique property.
(i) The calculus of $X \otimes Y$ equals its probabilistic coherence
space, and the calculus of $(X \otimes Y)^{\perp}$ equals its
probabilistic coherence space iff the strong product $G_X \boxtimes G_Y$
is perfect.
(ii) If $G_X$ and $G_Y$ are cographs, this holds iff $G_X$ or $G_Y$ has no
induced $4$-cycle.
\end{theorem}

\noindent
If the calculus equals the probabilistic coherence space at $X$ and at
$Y$, it does so at $X \otimes Y$, $X \mathbin{\&} Y$, and $X \oplus Y$,
and $P((X \otimes Y)^{\perp})$ is cut out by the products of
$\{0,1\}$-valued elements of $P(X)$ and $P(Y)$. The Lean proof derives a
bipolar theorem for downward-closed sets of $\{0,1\}$-vectors from
Theorem~\ref{thm:mix-eq-laws}. With the clique property these products
are the indicators of the products $K \times L$ of cliques, which are the
maximal cliques of $F = G_X \boxtimes G_Y$. So $P((X \otimes Y)^{\perp})
= \mathrm{QSTAB}(F)$, its $\{0,1\}$-valued elements are the stable sets of
$F$, and the calculus is $\mathrm{STAB}(F)$, which equals
$\mathrm{QSTAB}(F)$ iff $F$ is perfect~\cite{lovasz1972,chvatal1975}. In
(ii), if both graphs have an induced $4$-cycle the gap is
Theorem~\ref{thm:c4}. The closure statements, the description of
$P((X \otimes Y)^{\perp})$, and (ii) are formalized, and the equivalence
with perfection in (i) is cited.

\paragraph{A direct proof of (ii).}
The converse in (ii) has a direct proof, by an explicit decomposition,
that does not use perfect graph theory. Its Lean proof uses classical
case distinctions on real numbers, through the library. If neither side of $A \mathbin{\&} B$ is a clique, an
incoherent pair in $A$ and one in $B$ form an induced $4$-cycle. So a
formula built from $\mathbf{1}$ with $\mathbin{\&}$ and $\oplus$ that has
no induced $4$-cycle is isomorphic to one built from $\mathbf{1}$ with
$\oplus$ and $A \mapsto A \mathbin{\&} \mathbf{1}$, by a bijection of the
webs that preserves coherence in both directions\leanref{tree-of-noc4}.
Such bijections, and the exchange of the factors of $\otimes$, are
coordinate retracts, so we may assume that $X$ is built in this way. Its
points are the nodes of a rooted forest, the point of $\mathbf{1}$ in
$A \mathbin{\&} \mathbf{1}$ being the root above the forest of $A$, and
two points cohere iff one lies on the path from a root to the other. So
$G_X$ is the comparability graph of the forest (two nodes adjacent iff
one is an ancestor of the other)~\cite{wolk1962,golumbic1978}. For
$v : |Y| \to [0, \infty)$ let $\omega(v)$ be the largest weight of a
clique of $G_Y$, a sum at $\mathbin{\&}$ and a maximum at $\oplus$. With
$z_x = z(x, \cdot)$, a nonnegative $z$ on $|X| \times |Y|$ lies in
$P((X \otimes Y)^{\perp})$ iff $\omega(\sum_{x \in \pi} z_x) \le 1$ for
every path $\pi$ from a root\leanref{mem-p-iff}.

Every such $z$ is a mixture of $\{0,1\}$-valued functions whose sums
along the paths are indicators of stable sets of $G_Y$, and these are
elements of the calculus\leanref{decompose}. The proof is an induction on
the forest. At a root $r$ above a forest $X'$, the pair of $m = z_r$ and
the restriction $y$ of $z$ to $X'$ is a mixture of pairs $(1_I, y')$ with
$I$ stable and $\omega(1_I + \sum_{x \in \pi} y'_x) \le 1$ for every path
$\pi$ of $X'$, and the induction continues at $y'$. This root step is an
induction on $Y$. At $Y_1 \oplus Y_2$ the two components are decomposed
separately and combined with product weights. At $Y_1 \mathbin{\&} Y_2$
a stable set lies in one component. With $m_j$ the restriction of $m$ to
$Y_j$, the set $I$ is drawn in $Y_j$ with probability $\omega(m_j)$ and is
empty with probability $1 - \omega(m_1) - \omega(m_2)$. Before the draw,
the restriction $y_j$ of $y$ to $Y_j$ is split as $p_j + q_j$ with
$\omega(m_j + \sum_{x \in \pi} p_{j,x}) = \omega(m_j)$ and
$\omega(\sum_{x \in \pi} q_{j,x}) \le \omega(m_j + \sum_{x \in \pi}
y_{j,x}) - \omega(m_j)$ along every path $\pi$. The part $p_j$ is
assigned to the draws in $Y_j$ and $q_1 + q_2$ to the empty set. The
split is built by induction on $Y_j$ and a pass down the forest that
keeps the value at a node while $\omega$ stays within the level and, at
the first node where it does not, keeps the part that raises $\omega$
exactly to the level\leanref{absorb}.

Lean checks (ii) in both directions for $X$ and $Y$ built from
$\mathbf{1}$ with $\mathbin{\&}$ and $\oplus$\leanref{second-iff}. For
the other formulas in (ii) with nonempty webs, the clique property makes
a coherence-preserving bijection onto such a formula a coordinate retract
in both directions, and this transfer is not formalized. With (i), the
proof shows that the strong product of a forest comparability graph and
a cograph is perfect. Ravindra proved that the strong product of two
graphs with no induced path on four vertices and no induced $4$-cycle,
which are exactly the forest comparability graphs, is perfect, and
characterized the perfect strong products of two bipartite
graphs~\cite{ravindra1978}. The case of a forest comparability graph and
an arbitrary cograph extends the first result, and we did not find it in
the literature. For a single coherence space, an exercise of Ehrhard's
asks for the self-duality of its clique-constrained probabilistic
coherence space when the space is generated from $\mathbf{1}$ by
$\perp$ and $\mathbin{\&}$, so that its graph is a cograph, and for its
failure at the five-cycle~\cite{ehrhard-td2}. A proof through perfection is also
possible, and is not formalized: the product has no induced cycle of
length at least $5$ and no complement of one of length at least $6$, so
it is perfect by Hayward's theorem~\cite{hayward1985}.

For $X = \mathbf{n} \multimap \mathbf{m}$ the coherence graph is complete
multipartite (vertices in classes, adjacent iff in different classes)
with $n$ classes of size $m$, which has an induced $4$-cycle
iff $n, m \ge 2$. So the calculus of $T$ equals $T$ iff one of $X, Y, X',
Y'$ has at most one element, and at the core $(X \multimap Y) \otimes
(X' \multimap Y')$ itself the calculus equals the probabilistic coherence
space.

\paragraph{Beyond second order.}
The condition of Theorem~\ref{thm:c4} is not necessary in general.

\begin{proposition}[a gap without four-cycles\leanref{gap-f3}]
\label{prop:f3}
The calculus of $(\mathbf{2} \mathbin{\&} \mathbf{1}) \multimap
((\mathbf{2} \multimap \mathbf{2}) \otimes (\mathbf{2} \multimap
\mathbf{2}))$ is strictly smaller than its probabilistic coherence space.
Every tensor occurring negatively in it has a factor with at most three
points, so none has an induced $4$-cycle.
\end{proposition}

\noindent
The formula is $((\mathbf{2} \mathbin{\&} \mathbf{1}) \otimes T)^{\perp}$.
The witness has value $\tfrac12$ at five points of the web. Its membership
uses four incoherent pairs of $T$. The five points form a $5$-cycle of the
coherence graph, so every $\{0,1\}$-valued element is $1$ at no more than
two of them, and every mixture has total weight at most $2$ there,
against $\tfrac52$. The coherence graph of $T$ contains an induced house,
a $5$-cycle with one chord, and on fifteen points of the web the space is
$\mathrm{QSTAB}$ of the strong product of a path on three vertices with
the house, which is not perfect. Beyond second order the coherence graphs
are no longer cographs, and an exact condition would have to account for
the imperfect strong products of the graphs that arise.

The \emph{independence number} $\alpha(G)$ is the largest size of a
stable set, and the \emph{fractional packing number} $\alpha^{*}(G)$ is
the largest value of $\sum_v x_v$ over $\mathrm{QSTAB}(G)$. The bounds
$2$ and $\tfrac52$ of Theorem~\ref{thm:pentagon} are $\alpha(C_5)$ and
$\alpha^{*}(C_5)$. Cabello, Severini and Winter~\cite{cabello2014} show
that in a contextuality scenario (a set of outcomes with a family of
jointly measurable subsets) with exclusivity graph $G$ the noncontextual
bound is $\alpha(G)$, the quantum bound is the Lov\'asz number
$\vartheta(G)$ (the largest total weight on the theta body of $G$, a
convex set defined by orthonormal representations of the
graph)~\cite{lovasz1979}, and the bound for theories satisfying only
exclusivity is $\alpha^{*}(G)$. For $C_5$ this is the KCBS
inequality (the five-outcome noncontextuality inequality of Klyachko,
Can, Binicio\u{g}lu, and Shumovsky)~\cite{klyachko2008}, with
$\vartheta(C_5) = \sqrt5$. A calculus for $T$ corresponding to the theta
body of the exclusivity graph~\cite{gls1988}, which lies between $\mathrm{STAB}$ and
$\mathrm{QSTAB}$, is not constructed here. Section~\ref{sec:quantum}
makes the same comparison for a Kochen--Specker set.
